\exercice
\begin{em}
  Toutes les questions sont indépendantes.
\end{em}

On considère la fonction $((f))$ définie sur $\mathbb{R}$ par :
(* if nature ==  "degre1" *)
$(( f ))(x) = (( a|facteur("*x") ))(( b|facteur("so") ))$.
(* elif nature ==  "degre2" *)
$(( f ))(x) = (( a|facteur("*X") ))(( b|facteur("*sox") ))(( c|facteur("so") ))$.
(* endif *)

\begin{enumerate}
  (* if sens ==  "antecedent-vers-image" *)
  \item Quelles sont les coordonnées du point d'intersection entre la courbe de $(( f ))$ avec l'axe des ordonnées ?
  (* else *)
  \item Quelles sont les coordonnées des éventuels points d'intersection entre la courbe de $(( f ))$ et l'axe des abscisses ?
  (* endif *)

  \item Le point $A( (( A.x )); (( A.y )) )$ est-il sur la courbe de $(( f ))$ ?

  (* if sens ==  "antecedent-vers-image" *)
  \item Le point $B( (( B.x ));y)$ est sur la courbe de $(( f ))$. Quelle est la valeur de $y$ ?
  (* else *)
  \item Le point $B(x; (( B.y )))$ est sur la courbe de $(( f ))$. Quelles sont les valeurs possibles de $x$ ?
  (* endif *)
\end{enumerate}
